\documentclass[11pt,a4paper]{article}

\usepackage[utf8]{inputenc}
\usepackage[margin=1in]{geometry}
\usepackage{amsmath,amssymb,amsfonts,bm}
\usepackage{booktabs}
\usepackage{graphicx}
\usepackage{natbib}
\usepackage{hyperref}
\usepackage{microtype}
\usepackage{setspace}
\usepackage{caption}
\usepackage{subcaption}
\usepackage{tabularx}
\usepackage{array}

\hypersetup{
    colorlinks=true,
    linkcolor=blue,
    citecolor=blue,
    urlcolor=blue
}

\onehalfspacing

\title{\textbf{Economics-Aware Causal Marketing Targeting: Robust Policy Learning for Incremental Revenue Under Budget and Campaign-Cost Uncertainty}}

\author{
    \textbf{Department of Computer Science and Engineering} \\
    \textbf{\& Department of Business Administration} \\
    BRAC University, Dhaka, Bangladesh
}

\date{\today}

\begin{document}

\maketitle

\begin{abstract}
Traditional direct marketing analytics relies predominantly on predictive response models that estimate a customer's likelihood of purchase ($P(Y=1|X)$). However, targeting customers with high response probabilities frequently wastes marketing capital on ``Sure Things''---customers who would transact regardless of campaign exposure. Uplift modeling addresses this by estimating the Conditional Average Treatment Effect (CATE), shifting the analytical objective from prediction to causal incrementality. Despite algorithmic progress in causal machine learning, a critical decision gap persists: standard uplift formulations evaluate binary conversion lift without accounting for variable customer spend, unit campaign contact costs, product contribution margins, capacity constraints, or macroeconomic parameter uncertainty.

In this study, we introduce an integrated, economics-aware causal policy learning framework that maps heterogeneous treatment effects directly into budget-constrained commercial decisions. Using the benchmark Kevin Hillstrom \textit{MineThatData} randomized controlled trial ($N = 64,000$) and an external validation on the unbiased Criteo Uplift benchmark ($N \approx 13.98 \times 10^6$), we evaluate a taxonomy of response models, causal meta-learners ($S$-, $T$-, and $X$-Learners), doubly robust augmented inverse probability weighting (AIPW) estimators, and two-stage hurdle revenue uplift models. We formulate incremental customer economic utility ($U_i = \Delta R_i \cdot M - c$), individual break-even contact thresholds, budget-constrained knapsack policy allocation, and a minimax regret decision rule across a $4 \times 4$ commercial scenario grid ($M \in [0.20, 0.80]$, $c \in [\$0.10, \$1.00]$).

Empirical results reveal: (1) customer rankings produced by conventional response models diverge sharply from causal uplift models, exhibiting a Spearman rank correlation of only $\rho = 0.3817$ and merely $43.5\%$ cohort overlap within the top decile, confirming that response models misallocate over half their budget to non-incremental buyers; (2) at a $20\%$ campaign coverage constraint, the causal uplift policy generates $\$0.3934$ in net economic value per customer compared to $\$0.2908$ under response modeling---a $+35.3\%$ net profit expansion; (3) across uncertain economic scenarios, the proposed minimax regret policy compresses maximum financial regret by $61.1\%$ (from $\$0.3554$ to $\$0.1384$) relative to response targeting; (4) in multi-arm campaign allocation (Men's vs. Women's email vs. Control), uniform deployment of the globally superior creative outperforms individual algorithmic personalization ($\$0.3575$ vs. $\$0.2615$) due to ``opportunity cost drag''; and (5) large-scale digital display advertising (Criteo) exhibits high response-uplift concordance ($\rho = 0.8044$, $85.2\%$ overlap), establishing a fundamental domain boundary condition between high-intent direct channels and low-baseline programmatic display advertising.
\end{abstract}

\noindent \textbf{Keywords:} Causal Machine Learning, Uplift Modeling, Economic Decision Theory, Budget-Constrained Targeting, Minimax Regret, Marketing Analytics.

\newpage

\section{Introduction}

Direct marketing campaigns represent a major operational expenditure for modern retail, e-commerce, and direct-to-consumer (DTC) enterprises. Historically, marketing analysts have addressed customer targeting through the lens of supervised machine learning and predictive response modeling:
\begin{equation}
X_i \longrightarrow \widehat{P}(Y_i=1 \mid X_i) \longrightarrow \text{Target top deciles}
\end{equation}
where $X_i \in \mathbb{R}^d$ denotes observed pre-campaign customer characteristics (e.g., recency, purchase frequency, monetary history, channel preferences), and $Y_i \in \{0, 1\}$ represents the post-campaign binary transaction outcome \citep{blattberg2008database}. 

While computationally straightforward, this predictive paradigm suffers from a fundamental conceptual flaw: \textbf{high response probability does not imply campaign causality}. A customer with high purchase propensity may simply be an organic buyer whose transaction was imminent irrespective of receiving promotional outreach. When marketing budgets are allocated to these ``Sure Things'', promotional expenditures generate zero incremental revenue and simply dilute profit margins \citep{radcliffe2011real, devriendt2018literature}.

Uplift modeling resolves this limitation by operating within the Neyman-Rubin Potential Outcomes Framework \citep{rubin1974estimating}. Rather than predicting the conditional probability of purchase, uplift models estimate the Conditional Average Treatment Effect (CATE):
\begin{equation}
\tau(X_i) = \mathbb{E}[Y_i(1) - Y_i(0) \mid X_i]
\end{equation}
where $Y_i(1)$ and $Y_i(0)$ represent potential outcomes under promotional contact ($W_i=1$) and control ($W_i=0$), respectively. By identifying ``Persuadables''---customers whose transaction is strictly contingent upon intervention---causal targeting promises substantial efficiency gains \citep{gutierrez2017causal, kunzel2019metalearners}.

Despite widespread algorithmic advancements in causal machine learning, a significant divide persists between econometric CATE estimation and commercial decision-making in practice:
\begin{enumerate}
    \item \textbf{Conversion vs. Revenue Mismatch:} Most uplift algorithms focus strictly on binary conversion ($\Delta P$). However, customer monetary spend is highly zero-inflated and right-skewed. A customer with a modest conversion lift who purchases high-ticket items may generate substantially more incremental gross profit than a customer with a high conversion lift who buys discounted, low-margin merchandise \citep{lo2002true, gubela2019conversion}.
    \item \textbf{Neglect of Campaign Economics:} Causal models typically produce unitless rank scores or conversion probabilities. They do not incorporate product contribution margin rates ($M \in (0, 1]$) or unit contact/fulfillment costs ($c > 0$) directly into the objective function, obscuring the customer-level break-even threshold:
    \begin{equation}
    c^*_{i} = \Delta R_i \cdot M
    \end{equation}
    \item \textbf{Rigid Budget and Capacity Constraints:} Marketing managers rarely have unlimited capital. Campaigns operate under strict budget limits ($B$) or operational capacity caps ($K$). The allocation of multi-touch, multi-creative campaigns under fixed resources represents a constrained resource allocation problem rather than a simple unconstrained ranking exercise \citep{hauser2009non, sun2021budget}.
    \item \textbf{Economic Parameter Uncertainty:} In practice, marketers rarely know exact downstream contribution margins (due to returns, promotional markdowns, and inventory costs) or exact unit delivery fees. Policies optimized for a single assumed margin or cost structure risk catastrophic profit degradation if macroeconomic conditions shift. Decision frameworks must incorporate robustness against parameter ambiguity \citep{savage1951theory, manski2000identification}.
\end{enumerate}

To bridge this gap, this paper introduces an integrated, economics-aware causal policy learning framework and empirically validates it on two randomized controlled trials comprising over $14$ million customer observations.

\section{Theoretical Foundations \& Related Literature}

\subsection{Response Modeling vs. Causal Uplift}
In classical database marketing, the Recency, Frequency, and Monetary (RFM) paradigm has dominated customer selection \citep{bult1995optimal, fader2005rfm}. Customers with high past activity receive high response scores $P(Y=1|X)$ \citep{hughes1994strategic, verhoef2003commercial}. However, \citet{radcliffe2011real} formalized the four customer behavioral archetypes:
\begin{itemize}
    \item \textbf{Persuadables:} $Y(1) = 1, Y(0) = 0$ (Target group).
    \item \textbf{Sure Things:} $Y(1) = 1, Y(0) = 1$ (Wasteful spend).
    \item \textbf{Lost Causes:} $Y(1) = 0, Y(0) = 0$ (Unresponsive).
    \item \textbf{Sleeping Dogs:} $Y(1) = 0, Y(0) = 1$ (Negative impact).
\end{itemize}
Response models conflate Sure Things with Persuadables because both have high $P(Y=1|X)$. Uplift models isolate Persuadables by directly targeting $\tau(X)$ \citep{kane2014mining}.

\subsection{Causal Meta-Learners \& Orthogonal Scores}
\citet{kunzel2019metalearners} categorized generic meta-learners: the single-model $S$-Learner, twin-model $T$-Learner, and crossover $X$-Learner. For debiased estimation under high-dimensional covariates, \citet{chernozhukov2018double} developed Double Machine Learning (DML) based on Neyman-orthogonal score functions and cross-fitting, building on augmented inverse probability weighting \citep{robins1994estimation}. Tree-based causal forests \citep{athey2016recursive, wager2018estimation} optimize heterogeneity directly.

\subsection{Revenue Modeling \& Two-Stage Hurdle}
Because customer spend is zero-inflated, linear regressors produce unstable estimates. Econometric literature relies on two-part hurdle formulations \citep{cragg1971some, duan1983comparison, gubela2019conversion}:
\begin{equation}
\mathbb{E}[R \mid X, W] = P(Y=1 \mid X, W) \times \mathbb{E}[R \mid X, W, Y=1]
\end{equation}

\subsection{Robust Decision Theory \& Minimax Regret}
Under severe parameter ambiguity regarding contribution margins and contact costs, \citet{savage1951theory} and \citet{manski2000identification} advocate minimizing maximum regret:
\begin{equation}
\text{Regret}(\pi; \theta) = V(\pi^*_\theta; \theta) - V(\pi; \theta)
\end{equation}
which balances opportunity losses without the excessive conservatism of Wald's maximin criterion \citep{wald1945statistical, stoye2009minimax, bertsimas2016robust}.

\section{Economic Decision Framework}

\subsection{Incremental Economic Utility \& Break-Even Cost}
For customer $i$ and marketing action $a \in \{0, \dots, A\}$:
\begin{equation}
U_i(a) = \Delta R_{i, a} \cdot M - c_a
\end{equation}
where $\Delta R_{i, a} = \mathbb{E}[R_i(a) - R_i(0) \mid X_i]$ is the incremental revenue, $M$ is the contribution margin rate, and $c_a$ is the unit delivery cost ($c_0 = 0$).

The individual break-even contact threshold is:
\begin{equation}
c^*_{i, a} = \Delta R_{i, a} \cdot M
\end{equation}
Targeting customer $i$ with action $a$ is value-accretive if and only if $c_a \le c^*_{i, a}$.

\subsection{Budget-Constrained Optimization}
Given campaign budget $B$, the firm solves:
\begin{equation}
\max_{\pi} \sum_{i=1}^N U_i(\pi(X_i)) \quad \text{s.t.} \quad \sum_{i=1}^N c_{\pi(X_i)} \le B
\end{equation}
Under homogeneous costs ($c_a = c$), this simplifies to top-$K$ selection where $K = \lfloor B/c \rfloor$.

\subsection{Minimax Regret Policy}
Across operating scenarios $\theta = (M, c) \in \Theta$:
\begin{equation}
\pi_R = \arg\min_{\pi \in \Pi} \max_{\theta \in \Theta} \left[ V(\pi^*_\theta; \theta) - V(\pi; \theta) \right]
\end{equation}

\section{Data and Covariate Balance Audit}

\subsection{Kevin Hillstrom MineThatData RCT}
The Kevin Hillstrom dataset \citep{hillstrom2008minethatdata} contains $64,000$ retail customers randomly assigned to Men's E-Mail ($N = 21,307$), Women's E-Mail ($N = 21,387$), or Control ($N = 21,306$). 

Table~\ref{tab:balance} confirms rigorous pre-treatment covariate balance across all eight observed features. The maximum Standardized Mean Difference (SMD) is $0.0169$ ($< 0.05$ threshold), and an omnibus 5-fold cross-validated logistic classifier achieves an ROC-AUC of $0.4991$ ($\approx 0.5000$), confirming perfect randomization.

\begin{table}[t]
\centering
\small
\caption{Pre-Treatment Covariate Balance Audit (Standardized Mean Differences)}
\label{tab:balance}
\begin{tabular}{lcccccc}
\toprule
Feature & Mean Ctrl & Mean Men's & SMD (Men's) & Mean Women's & SMD (Women's) & Max SMD \\
\midrule
\texttt{recency} & 5.750 & 5.774 & 0.0068 & 5.768 & 0.0052 & 0.0068 \\
\texttt{history} & 240.88 & 242.84 & 0.0076 & 242.54 & 0.0065 & 0.0076 \\
\texttt{mens} & 0.553 & 0.551 & 0.0046 & 0.549 & 0.0086 & 0.0086 \\
\texttt{womens} & 0.548 & 0.551 & 0.0076 & 0.550 & 0.0049 & 0.0076 \\
\texttt{newbie} & 0.502 & 0.502 & 0.0009 & 0.503 & 0.0026 & 0.0034 \\
\texttt{zip\_Suburban} & 0.452 & 0.446 & 0.0117 & 0.451 & 0.0011 & 0.0117 \\
\texttt{zip\_Urban} & 0.401 & 0.402 & 0.0020 & 0.400 & 0.0018 & 0.0037 \\
\texttt{channel\_Phone} & 0.438 & 0.434 & 0.0083 & 0.442 & 0.0086 & 0.0169 \\
\texttt{channel\_Web} & 0.440 & 0.445 & 0.0110 & 0.437 & 0.0051 & 0.0162 \\
\bottomrule
\end{tabular}
\end{table}

\subsection{Criteo Uplift Benchmark}
The Criteo dataset \citep{diemert2018large} contains $13,979,592$ impressions from randomized ad auctions ($85.00\%$ treatment share, baseline conversion $0.2917\%$).

\section{Empirical Results}

\subsection{Hypothesis H1: Response vs. Uplift Ranking Divergence}
Evaluating Model A2 (XGBoost Response) against Model B2 ($T$-Learner Conversion Uplift) on held-out test data ($N = 8,523$) yields a Spearman rank correlation of $\rho = 0.3817$ ($p = 9.70 \times 10^{-294}$). Within the top $10\%$ targeting cohort, the overlap is only $43.5\%$. This confirms Hypothesis H1: traditional response models allocate $56.5\%$ of their marketing resources to non-incremental buyers.

\begin{table}[h]
\centering
\small
\caption{Model Qini Scores and Agreement with Response Scoring}
\label{tab:qini}
\begin{tabular}{llcc}
\toprule
Model ID & Model Name & Qini Score (Conv) & Spearman $\rho$ vs. Response \\
\midrule
Model A2 & XGBoost Response & -2.1997 & 1.0000 \\
Model B1 & $S$-Learner & -4.1986 & 0.3124 \\
Model B2 & $T$-Learner & \textbf{-1.7132} & \textbf{0.3817} \\
Model B3 & $X$-Learner & -3.7856 & 0.3450 \\
Model C & Doubly Robust AIPW & -4.3984 & 0.2981 \\
Hurdle & Two-Stage Hurdle & -0.9369 & 0.3683 \\
\bottomrule
\end{tabular}
\end{table}

\subsection{Hypotheses H2 \& H3: Economic Value Creation}
Table~\ref{tab:policy} reports held-out policy values at $20\%$ coverage ($M = 0.50, c = \$0.05$). The $T$-Learner uplift policy achieves a net economic value of $\$0.3934$ per customer, delivering a $+35.28\%$ net profit increase over the Response policy ($\$0.2908$).

\begin{table}[h]
\centering
\small
\caption{Held-Out Policy Performance ($20\%$ Coverage Constraint)}
\label{tab:policy}
\begin{tabular}{lcccc}
\toprule
Policy & Targeted (\%) & Exp. Conv. (\%) & Exp. Spend (\$) & Net Value (\$/cust) \\
\midrule
Random Policy & 20.0\% & 0.6805\% & \$0.6812 & \$0.2906 \\
Response Policy (XGB) & 20.0\% & 0.5163\% & \$0.6817 & \$0.2908 \\
$T$-Learner Uplift Policy & 20.0\% & \textbf{0.6336\%} & \textbf{\$0.8869} & \textbf{\$0.3934} \\
Revenue Uplift Policy & 20.0\% & 0.5866\% & \$0.8160 & \$0.3580 \\
Profit-Aware Policy & 20.0\% & 0.5866\% & \$0.8160 & \$0.3580 \\
\bottomrule
\end{tabular}
\end{table}

\subsection{Hypotheses H4 \& H5: Minimax Regret Robustness}
Across 16 economic operating scenarios ($M \in \{0.20, 0.40, 0.60, 0.80\}$, $c \in \{\$0.10, \$0.25, \$0.50, \$1.00\}$), the conventional Response policy incurs a maximum economic regret of $\$0.3554$ per customer. The Profit-Aware Robust Policy caps maximum regret at $\$0.1384$---a $61.06\%$ reduction in maximum downside financial regret.

\begin{table}[h]
\centering
\small
\caption{Policy Regret Summary Across 16 Economic Scenarios}
\label{tab:regret}
\begin{tabular}{lcccc}
\toprule
Metric & Response Policy & Uplift Policy & Revenue Uplift & Profit-Aware Robust \\
\midrule
Maximum Regret & \$0.3554 & \$0.1912 & \$0.2479 & \textbf{\$0.1384} \\
Mean Regret & \$0.1341 & \$0.0315 & \$0.0669 & \textbf{\$0.0260} \\
\bottomrule
\end{tabular}
\end{table}

\subsection{Hypothesis H6: Multi-Arm Personalization \& Opportunity Cost Drag}
In Experiment B ($3$ arms, $N = 12,800$ test customers), the Uniform Men's Email campaign generates $\$0.3575$ net value per customer, whereas the personalized causal policy generates $\$0.2615$ ($\Delta V = -\$0.0960$, $p = 0.2720$). Because Men's email has a universally superior Average Treatment Effect ($\text{ATE} = +\$0.77$ vs. $+\$0.42$ for Women's email), routing customers to Women's email based on minor sub-cohort preferences creates an \textit{opportunity cost drag}, eroding net returns.

\subsection{Hypothesis H7: Large-Scale Display Advertising Boundary Condition}
On Criteo's $13.98$-million impression dataset, Response XGBoost and $T$-Learner Uplift exhibit a Spearman rank correlation of $\rho = 0.8044$ and $85.2\%$ top-decile overlap. In programmatic display advertising with $85\%$ treatment saturation and low organic conversion ($0.29\%$), response scoring effectively proxies uplift. This establishes a key boundary condition: causal uplift modeling is critical for high-intent push channels (direct mail, email), whereas response models remain competitive in saturated display ad environments.

\section{Managerial Implications}

\begin{enumerate}
    \item \textbf{Transition to Break-Even Filtering:} Marketing executives must eliminate pure response-propensity ranking and enforce the break-even condition $c^*_i = \Delta R_i \cdot M$.
    \item \textbf{Optimal Campaign Capacity:} Policy value curves demonstrate that net profit peaks between $15\%$ and $25\%$ coverage, declining sharply beyond $30\%$ as campaigns penetrate un-persuadable cohorts.
    \item \textbf{Hedging Macroeconomic Volatility:} In volatile margin environments, deploying minimax regret policies reduces maximum profit downside by $61.1\%$ while preserving cohort stability.
    \item \textbf{Creative Dominance Check:} Before investing in micro-personalization infrastructure, firms must test for creative dominance. When one creative has a dominant ATE, uniform rollout is economically optimal.
\end{enumerate}

\section{Conclusion}
This study bridges the gap between econometric causal inference and commercial marketing decision-making. By integrating causal meta-learners, two-stage hurdle spend models, knapsack budgeting, and minimax regret optimization, we provide a robust decision architecture for profit-maximizing marketing interventions.

\bibliographystyle{apalike}
\bibliography{references}

\end{document}
